\documentclass[12pt]{article}

% Idioma y codificación
\usepackage[spanish, es-tabla]{babel}       %es-tabla para que se titule "Tabla"
\usepackage[utf8]{inputenc}

% Márgenes
\usepackage[a4paper,top=3cm,bottom=2.5cm,left=3cm,right=3cm]{geometry}

% Comentarios de bloque
\usepackage{verbatim}

% Paquetes de links
\usepackage[hidelinks]{hyperref}    % Permite enlaces
\usepackage{url}                    % redirecciona a la web

% Más opciones para enumeraciones
\usepackage{enumitem}


%------------------------------------------------------------------------

%Paquetes de figuras
\usepackage{caption}
\usepackage{subcaption} % Figuras al lado de otras
\usepackage{float}      % Poner figuras en el sitio indicado H.


% Paquetes de imágenes
\usepackage{graphicx}       % Paquete para añadir imágenes
\usepackage{transparent}    % Para manejar la opacidad de las figuras

% Paquete para usar colores
\usepackage[dvipsnames]{xcolor}
\usepackage{pagecolor}      % Para cambiar el color de la página

% Habilita tamaños de fuente mayores
\usepackage{fix-cm}


%------------------------------------------------------------------------

% Paquetes de matemáticas
\usepackage{amsmath, amsfonts, amssymb, mathrsfs}
\usepackage{cancel}     % Simplificar tachando
\usepackage{polynom}    % Divisiones y Ruffini

\usepackage{pgfplots}   %Representar funciones
\pgfplotsset{compat=1.18}  % Versión 1.18

%Definición de teoremas, etc.
\usepackage{amsthm}
%\swapnumbers   % Intercambia la posición del texto y de la numeración

\theoremstyle{plain}
\newtheorem{teo}{Teorema}[chapter]     % Se resetea en cada chapter
\newtheorem{coro}{Corolario}[teo]           % Se resetea en cada teorema
\newtheorem{prop}[teo]{Proposición}         % Usa el mismo contador que teorema
\newtheorem{lema}[teo]{Lema}                % Usa el mismo contador que teorema

\theoremstyle{remark}
\newtheorem*{observacion}{Observación}

\theoremstyle{definition}
\newtheorem{definicion}{Definición} [chapter]     % Se resetea en cada chapter
\newtheorem*{notacion}{Notacion}
\newtheorem*{ejemplo}{Ejemplo}
\newtheorem*{ejercicio*}{Ejercicio}             % No numerado
\newtheorem{ejercicio}{Ejercicio} [section]     % Se resetea en cada section


% Modificar el formato de la numeración del teorema "ejercicio"
\renewcommand{\theejercicio}{%
  \ifnum\value{section}=0 % Si no se ha iniciado ninguna sección
    \arabic{ejercicio}% Solo mostrar el número de ejercicio
  \else
    \thesection.\arabic{ejercicio}% Mostrar número de sección y número de ejercicio
  \fi
}


% \renewcommand\qedsymbol{$\blacksquare$}         % Cambiar símbolo QED
%------------------------------------------------------------------------

% Paquetes para encabezados
\usepackage{fancyhdr}
\pagestyle{fancy}
\fancyhf{}

\newcommand{\helv}{ % Modificación tamaño de letra
\fontfamily{}\fontsize{12}{12}\selectfont}
\setlength{\headheight}{15pt} % Amplía el tamaño del índice


%\usepackage{lastpage}   % Referenciar última pag   \pageref{LastPage}
\fancyfoot[C]{\thepage}

%------------------------------------------------------------------------

\begin{comment} % Descomentar en book. Comentar en article

% Conseguir que no ponga "Capítulo 1". Sino solo "1."
\makeatletter
\renewcommand{\@makechapterhead}[1]{%
  \vspace*{50\p@}%
  {\parindent \z@ \raggedright \normalfont
    \ifnum \c@secnumdepth >\m@ne
      \huge\bfseries \thechapter.\hspace{1em}\ignorespaces
    \fi
    \interlinepenalty\@M
    \Huge \bfseries #1\par\nobreak
    \vskip 40\p@
  }}
\makeatother
\end{comment}

%------------------------------------------------------------------------
% Personalizar la portada
\usepackage{titling}


%------------------------------------------------------------------------

% Comandos definidos
\newcommand{\bb}[1]{\mathbb{#1}}
\newcommand{\cc}[1]{\mathcal{#1}}
\newcommand{\Hop}{\stackrel{L'H\hat{o}pital}{=}}

\newcommand{\red}[1]{{\color{red}#1}} % Para integrales, destacar los cambios.

% Método de integración
\newcommand{\MetInt}[2]{
    \left[\begin{array}{c}
        #1 \\ #2
    \end{array}\right]
}

%------------------------------------------------------------------------
\newcommand{\T}[0]{\cc{T}}

\begin{document}

    % 1. Foto de fondo
    % 2. Título
    % 3. Encabezado Izquierdo
    % 4. Color de fondo
    % 5. Coord x del titulo
    % 6. Coord y del titulo
    % 7. Fecha

    
    % 1. Foto de fondo
% 2. Título
% 3. Encabezado Izquierdo
% 4. Color de fondo
% 5. Coord x del titulo
% 6. Coord y del titulo
% 7. Fecha

\newcommand{\portada}[7]{

    \portadaBase{#1}{#2}{#3}{#4}{#5}{#6}{#7}
    \portadaBook{#1}{#2}{#3}{#4}{#5}{#6}{#7}
}

\newcommand{\portadaExamen}[7]{

    \portadaBase{#1}{#2}{#3}{#4}{#5}{#6}{#7}
    \portadaArticle{#1}{#2}{#3}{#4}{#5}{#6}{#7}
}




\newcommand{\portadaBase}[7]{

    % Tiene la portada principal y la licencia Creative Commons
    
    % 1. Foto de fondo
    % 2. Título
    % 3. Encabezado Izquierdo
    % 4. Color de fondo
    % 5. Coord x del titulo
    % 6. Coord y del titulo
    % 7. Fecha
    
    
    \thispagestyle{empty}               % Sin encabezado ni pie de página
    \newgeometry{margin=0cm}        % Márgenes nulos para la primera página
    
    
    % Encabezado
    \fancyhead[L]{\helv #3}
    \fancyhead[R]{\helv \nouppercase{\leftmark}}
    
    
    \pagecolor{#4}        % Color de fondo para la portada
    
    \begin{figure}[p]
        \centering
        \transparent{0.3}           % Opacidad del 30% para la imagen
        
        \includegraphics[width=\paperwidth, keepaspectratio]{assets/#1}
    
        \begin{tikzpicture}[remember picture, overlay]
            \node[anchor=north west, text=white, opacity=1, font=\fontsize{60}{90}\selectfont\bfseries\sffamily, align=left] at (#5, #6) {#2};
            
            \node[anchor=south east, text=white, opacity=1, font=\fontsize{12}{18}\selectfont\sffamily, align=right] at (9.7, 3) {\textbf{\href{https://losdeldgiim.github.io/}{Los Del DGIIM}}};
            
            \node[anchor=south east, text=white, opacity=1, font=\fontsize{12}{15}\selectfont\sffamily, align=right] at (9.7, 1.8) {Doble Grado en Ingeniería Informática y Matemáticas\\Universidad de Granada};
        \end{tikzpicture}
    \end{figure}
    
    
    \restoregeometry        % Restaurar márgenes normales para las páginas subsiguientes
    \pagecolor{white}       % Restaurar el color de página
    
    
    \newpage
    \thispagestyle{empty}               % Sin encabezado ni pie de página
    \begin{tikzpicture}[remember picture, overlay]
        \node[anchor=south west, inner sep=3cm] at (current page.south west) {
            \begin{minipage}{0.5\paperwidth}
                \href{https://creativecommons.org/licenses/by-nc-nd/4.0/}{
                    \includegraphics[height=2cm]{assets/Licencia.png}
                }\vspace{1cm}\\
                Esta obra está bajo una
                \href{https://creativecommons.org/licenses/by-nc-nd/4.0/}{
                    Licencia Creative Commons Atribución-NoComercial-SinDerivadas 4.0 Internacional (CC BY-NC-ND 4.0).
                }\\
    
                Eres libre de compartir y redistribuir el contenido de esta obra en cualquier medio o formato, siempre y cuando des el crédito adecuado a los autores originales y no persigas fines comerciales. 
            \end{minipage}
        };
    \end{tikzpicture}
    
    
    
    % 1. Foto de fondo
    % 2. Título
    % 3. Encabezado Izquierdo
    % 4. Color de fondo
    % 5. Coord x del titulo
    % 6. Coord y del titulo
    % 7. Fecha


}


\newcommand{\portadaBook}[7]{

    % 1. Foto de fondo
    % 2. Título
    % 3. Encabezado Izquierdo
    % 4. Color de fondo
    % 5. Coord x del titulo
    % 6. Coord y del titulo
    % 7. Fecha

    % Personaliza el formato del título
    \pretitle{\begin{center}\bfseries\fontsize{42}{56}\selectfont}
    \posttitle{\par\end{center}\vspace{2em}}
    
    % Personaliza el formato del autor
    \preauthor{\begin{center}\Large}
    \postauthor{\par\end{center}\vfill}
    
    % Personaliza el formato de la fecha
    \predate{\begin{center}\huge}
    \postdate{\par\end{center}\vspace{2em}}
    
    \title{#2}
    \author{\href{https://losdeldgiim.github.io/}{Los Del DGIIM}}
    \date{Granada, #7}
    \maketitle
    
    \tableofcontents
}




\newcommand{\portadaArticle}[7]{

    % 1. Foto de fondo
    % 2. Título
    % 3. Encabezado Izquierdo
    % 4. Color de fondo
    % 5. Coord x del titulo
    % 6. Coord y del titulo
    % 7. Fecha

    % Personaliza el formato del título
    \pretitle{\begin{center}\bfseries\fontsize{42}{56}\selectfont}
    \posttitle{\par\end{center}\vspace{2em}}
    
    % Personaliza el formato del autor
    \preauthor{\begin{center}\Large}
    \postauthor{\par\end{center}\vspace{3em}}
    
    % Personaliza el formato de la fecha
    \predate{\begin{center}\huge}
    \postdate{\par\end{center}\vspace{5em}}
    
    \title{#2}
    \author{\href{https://losdeldgiim.github.io/}{Los Del DGIIM}}
    \date{Granada, #7}
    \thispagestyle{empty}               % Sin encabezado ni pie de página
    \maketitle
    \vfill
}
    \portadaExamen{ffccA4.jpg}{Topología I\\Examen I}{Topología I. Examen I}{MidnightBlue}{-8}{28}{2023-2024}{Arturo Olivares Martos}

    \begin{description}
        \item[Asignatura] Topología I.
        \item[Curso Académico] 2023-24.
        \item[Grado] Doble Grado en Ingeniería Informática y Matemáticas.
        \item[Grupo] Único.
        \item[Profesor] Antonio Alarcón López.
        \item[Descripción] Parcial del Tema 1.
        \item[Fecha] 10 de noviembre de 2023.
        \item[Duración] 60 minutos.
    
    \end{description}
    \newpage
    
    \begin{ejercicio}[4 puntos]
        Demuestra que el primer axioma de separación y el primer axioma de numerabilidad son hereditarios.
    \end{ejercicio}

    \begin{ejercicio}[6 puntos]
        En el conjunto $\bb{N}=\{1,2,\dots\}$  de los números naturales, se consideran las siguientes topologías:
        \begin{equation*}
            \begin{split}
                \T_1&=\{\emptyset, \bb{N}\}\cup \{[n,+\infty[~\cap \bb{N}\mid n\in \bb{N}\}\\\T_2&=\{\emptyset, \bb{N}\}\cup \{[1,n]\cap \bb{N}\mid n\in \bb{N}\}
            \end{split}
        \end{equation*}
        \begin{enumerate}
            \item Para cada $n\in \bb{N}$, determinar la base de entornos de $n$ más económica posible (es decir, la que tenga menor cantidad de entornos posible), en ambas topologías.

            Dado $n\in \bb{N}$, veamos que una base de entornos de $n$ en $\T_1$ es:
            \begin{equation*}
                \beta_n^1=\{[n,+\infty[~\cap \bb{N}\} = \{ \{ n, n+1, \dots \} \}
            \end{equation*}
            En primer lugar, es directo ver que como $n\in [n,+\infty[~\cap \bb{N}\in \T_1$, entonces $\beta_n^1$ está formada por un entorno de $n$.
            Sea ahora otro entorno $N$ de $n$. Entonces, $\exists m\in \bb{N}$ tal que $n\in [m,+\infty[~\cap \bb{N}\subset N$.
            Como $n\in [m,+\infty[$, entonces $n\geq m$, por lo que $[n,+\infty[~\subset [m,+\infty[$, y por tanto $[n,+\infty[~\cap \bb{N}\subset [m,+\infty[~\cap \bb{N}\subset N$.\\

            Veamos ahora que una base de entornos de $n$ en $\T_2$ es:
            \begin{equation*}
                \beta_n^2=\{[1,n]\cap \bb{N}\} = \{ \{ 1, 2, \dots, n \} \}
            \end{equation*}
            En primer lugar, es directo ver que como $n\in [1,n]\cap \bb{N}\in \T_2$, entonces $\beta_n^2$ está formada por un entorno de $n$.
            Sea ahora otro entorno $N$ de $n$. Entonces, $\exists m\in \bb{N}$ tal que $n\in [1,m]\cap \bb{N}\subset N$.
            Como $n\in [1,m]$, entonces $n\leq m$, por lo que $[1,n]\subset [1,m]$, y por tanto $[1,n]\cap \bb{N}\subset [1,m]\cap \bb{N}\subset N$.

            \item Calcula el interior, la adherencia y la frontera del conjunto $B=\{1,3,5\}$ en ambas topologías.
            
            Calculamos primero el interior de $B$ en $\T_1$. En primer lugar, sabemos que $\emptyset\subset B^\circ\subset B$. Además, como $B^\circ\in \T_1$, 
            y los abiertos de dicha topología son $\emptyset$ y conjuntos no acotados superiormente, entonces $B^\circ=\emptyset$.

            Calculamos ahora la adherencia de $B$ en $\T_1$. En primer lugar, sabemos que $B\subset \ol{B}\subset \bb{N}$. Además, como $\ol{B}\in C_{\T_1}$, entonces
            $\ol{B}=[1,n]\cap \bb{N}$ para $n\geq 5$. Como $\ol{B}$ es el menor cerrado que contiene a $B$, entonces $\ol{B}=[1,5]\cap \bb{N}=\{1,2,3,4,5\}$.

            Por tanto, en $\T_1$ tenemos que:
            \begin{equation*}
                \partial B = \ol{B}\setminus B^\circ = \{1,2,3,4,5\}\setminus \emptyset = \{1,2,3,4,5\}
            \end{equation*}

            Calculamos ahora el interior de $B$ en $\T_2$. En primer lugar, sabemos que por las propiedades del interior $\emptyset\subset B^\circ\subset B$. Además, como $B^\circ\in \T_2$,
            y los abiertos de dicha topología son $\emptyset$ y los conjuntos de números naturales sucesivos, entonces $B^\circ=\{1\}$.

            Calculamos ahora la adherencia de $B$ en $\T_2$. En primer lugar, sabemos que $B\subset \ol{B}\subset \bb{N}$. Además, como $\ol{B}\in C_{\T_2}$, entonces
            $\ol{B}=[n, +\infty[~\cap \bb{N}$ para $n\leq 1$. Como $n\in \bb{N}$, entonces $\ol{B}=\bb{N}$ (tenemos entonces que $B$ es denso). Por tanto, en $\T_2$ tenemos que:
            \begin{equation*}
                \partial B = \ol{B}\setminus B^\circ = \bb{N}\setminus \{1\}
            \end{equation*}

            \item Sea $A\subset \bb{N}$ el conjunto de los números impares.
            Calcula el interior, la adherencia, y la frontera del conjunto $B=\{1,3,5\}$ en $A$ con la topología inducida por cada una de las topologías.\\

            Tenemos que $B=[1,5]\cap A$. Por tanto, tenemos que $B\in C_{{\T_1}_{\big| A}}$, por lo que $\ol{B}^A = B$.
            Respecto al interior de $B$ en $A$ con la topología inducida por $\T_1$, sabemos que $\emptyset\subset B^{\circ_A}\subset B$. Además, como $B^{\circ_A}\in {\T_1}_{\big| A}$,
            y los abiertos de dicha topología son $\emptyset$ y conjuntos de impares no acotados superiormente, entonces $B^{\circ_A}=\emptyset$.
            
            Por tanto, en $A$ con la topología inducida por $\T_1$ tenemos que:
            \begin{equation*}
                \partial_A B = \ol{B}^A\setminus B^{\circ_A} = B\setminus \emptyset = B
            \end{equation*}

            Como $B=[1,5]\cap A$, entonces $B\in {\T_2}_{\big| A}$, por lo que $B^{\circ_A} = B$.
            Respecto a la adherencia de $B$ en $A$ con la topología inducida por $\T_2$, sabemos que $B\subset \ol{B}^A\subset A$.
            Además, como $\ol{B}^A\in C_{{\T_2}_{\big| A}}$, entonces $\ol{B}^A=[n,+\infty[~\cap A$ para $n\leq 1$. Como $n\in A$, entonces $\ol{B}^A=A$ (tenemos entonces que $B$ es denso).
            Por tanto, en $A$ con la topología inducida por $\T_2$ tenemos que:
            \begin{equation*}
                \partial_A B = \ol{B}^A\setminus B^{\circ_A} = A\setminus B = \{7, 9, 11, \dots\}
            \end{equation*}


            \item Determina si alguna de las topologías cumple alguno de los axiomas de separación.
            
            Veamos que ninguna de las topologías cumple el axioma T1, por lo que tampoco cumplirán el axioma T2. Para ello, consideramos los puntos $1$ y $2$.

            Para demostrar que $(\bb{N}, \T_1)$ no cumple el axioma T1, buscamos $U\in \T_1$ tal que $1\in U$ y $2\notin U$. Como $1\in U$, $U=\bb{N}$. Pero entonces $2\in U$, por lo que no existe tal $U$.

            Para demostrar que $(\bb{N}, \T_2)$ no cumple el axioma T1, buscamos $U\in \T_2$ tal que $2\in U$ y $1\notin U$. Como $2\in U$, $U=[1,n]\cap \bb{N}$ para algún $n\geq 2$.
            Pero entonces $1\in U$, por lo que no existe tal $U$.
        \end{enumerate}
    \end{ejercicio}
\end{document}